\section{换元与分部积分：反向识别求导规则}
\label{sec:02-Calculus/04-Integration/03}

拿到 \(\int 2x\cos(x^2)\,dx\) 的时候，你可能先翻了基本积分表。\(\int\cos u\,du=\sin u+C\)，没问题——但被积函数是 \(2x\cos(x^2)\)，不是 \(\cos u\) 搭上 \(du\)。表和眼前的东西对不上型号。

盯着它看几秒钟。\(\sin(x^2)\) 的导数是什么？

\[
\frac{d}{dx}\sin(x^2)=\cos(x^2)\cdot 2x=2x\cos(x^2).
\]

恰好就是被积函数。答案直接写出：\(\sin(x^2)+C\)。

这次凭眼力猜中了。换一个：\(\int x\cos(x^2)\,dx\)。还好——差一个常数因子 \(\frac12\)，凑一凑就行。再换：\(\int\cos(x^2)\,dx\)。这下没有前面的 \(x\) 帮忙了，不再是任何初等函数的导数。猜不出来了。

本节的两件工具——换元积分和分部积分——有一个共同的想法：把求导规则倒过来读。换元倒着读链式法则，分部倒着读乘积法则。它们不保证每次都积得出来——那是做不到的——但给了你一套比''盯着函数发呆等灵感''系统得多的搜索方法。

\subsection{换元：链式法则的逆向执行}

链式法则正向顺滑无比：

\[
\frac{d}{dx}F(g(x))=F'(g(x))\,g'(x).
\]

反过来读：如果你在被积式中同时发现了某个内层函数 \(g(x)\) 和它的导数因子 \(g'(x)\)，令 \(u=g(x)\) 就可以把积分翻转为

\[
\int F'(g(x))\,g'(x)\,dx = \int F'(u)\,du = F(u)+C = F(g(x))+C.
\]

操作步骤：设 \(u=g(x)\)，计算 \(du=g'(x)\,dx\)，然后把原积分中所有含 \(x\) 的因子连同那个 \(dx\) 一并替换成 \(u\) 和 \(du\)。

回到开头那个一眼看穿的例子：\(\int 2x\cos(x^2)\,dx\)。内层 \(x^2\)，它的导数 \(2x\) 正好在乘法里。设 \(u=x^2\)，\(du=2x\,dx\)。替换：\(\int\cos u\,du=\sin u+C=\sin(x^2)+C\)。不再靠眼力猜——靠的是识读''内层和导数因子同时在场''这个结构信号。

练一个需要补常数的：

\[
\int x\sqrt{x^2+1}\;dx.
\]

内层 \(x^2+1\)，导数 \(2x\)。原式有 \(x\,dx\) 但缺系数 2。设 \(u=x^2+1\)，\(du=2x\,dx\)，于是 \(x\,dx=\frac12 du\)：

\[
\int \sqrt{u}\cdot\frac12\,du = \frac12\cdot\frac{2}{3}u^{3/2}+C = \frac13(x^2+1)^{3/2}+C.
\]

再练一个定积分版：

\[
\int_0^{\sqrt{\pi}} 2x\cos(x^2)\,dx.
\]

设 \(u=x^2\)，\(du=2x\,dx\)。边界同步：\(x=0\Rightarrow u=0\)，\(x=\sqrt{\pi}\Rightarrow u=\pi\)。

\[
\int_0^{\pi} \cos u\,du = \bigl[\sin u\bigr]_0^{\pi} = \sin\pi-\sin0 = 0-0 = 0.
\]

图形上，\(\cos(x^2)\) 在 \([0,\sqrt{\pi}]\) 上的正负部分恰好抵消。没有几何直觉也要相信这个零——它是代数替换和基本定理联手得出的必然结果。

\subsection{换元的边界条件：什么情况下走不通}

换元要求内层函数的导数因子出现在被积式中——至少差一个常数。如果这个因子不存在，硬设 \(u=g(x)\) 只会制造一个更复杂的积分。

试试对 \(\int\cos(x^2)\,dx\) 设 \(u=x^2\)，\(du=2x\,dx\)。被积式里没有 \(x\,dx\)——你无法合法地把 \(dx\) 转换成只含 \(du\) 的表达式。不满足结构条件，换元帮不上忙。这不是你技术不够，是积分本身在初等函数范围内就没有闭合的原函数。

另一个常见坑：定积分换元忘记同步改边界。新被积变量配旧边界，或者换回原变量后代的却是新边界——这两种混搭都会产生无意义的数值。方法有两个等价版本：换元后直接用新边界计算；或者先求不定积分再换回原变量代原边界。只选一条路线走到底，不在中间换车道。

\(du=g'(x)\,dx\) 在一维积分里是个乘数修正。到了多重积分，同一思路会升级为 Jacobian 行列式，记录坐标变换在每个方向的局部面积或体积伸缩。此刻打下的''换元同时换小微元尺度''的直觉，以后直接平移。

\subsection{分部积分：乘积法则的逆向执行}

换元处理''内层加导数因子''的模式。但 \(\int x e^x\,dx\) 里，\(x\) 和 \(e^x\) 是独立的两块，没有内层关系。换元走不通。

回头看乘积法则：

\[
(uv)' = u'v + uv'.
\]

两边同时对 \(x\) 积分。\(uv\) 是 \((uv)'\) 的原函数，故

\[
uv = \int u'v\,dx + \int uv'\,dx = \int v\,du + \int u\,dv.
\]

移项重排得到分部积分公式：

\[
\int u\,dv = uv - \int v\,du.
\]

它把一个难积的乘积转成另一个积分——赌的是右边那个更简单。

用 \(\int x e^x\,dx\) 检验。令 \(u=x\)，\(dv=e^x\,dx\)，则 \(du=dx\)，\(v=e^x\)：

\[
\int x e^x\,dx = x e^x - \int e^x\,dx = x e^x - e^x + C = e^x(x-1)+C.
\]

被积函数 \(x e^x\) 本身无从下手，但拆成 \(u\) 和 \(dv\) 两块之后，新积分 \(\int e^x\,dx\) 直接就是 \(e^x\)。一趟分部就把一个看不穿的积分化成了一个基础积分。

\(\int x^2 e^x\,dx\) 可以看见一个模式：幂函数 \(x^2\) 和指数函数 \(e^x\) 的乘积。沿用同样的策略——令 \(u=x^2\)（求导降次），\(dv=e^x\,dx\)。则 \(du=2x\,dx\)，\(v=e^x\)：

\[
\int x^2 e^x\,dx = x^2 e^x - \int 2x e^x\,dx.
\]

右边的新积分 \(\int 2x e^x\,dx\) 比原来的次数低了一阶——它是我们刚才已经会算的 \(\int x e^x\,dx\) 的两倍。再分部一次：

\[
\int x^2 e^x\,dx = x^2 e^x - 2\bigl[e^x(x-1)\bigr] + C = e^x(x^2-2x+2)+C.
\]

这里有一个比结果更值得带走的结构观察：分部积分可以把幂函数一阶一阶''削平''。\(x^2\) 经一次分部降为 \(2x\)，两次分部后幂函数完全消失，剩下的只是指数项。这个模式——用分部削掉多项式因子——在 \(\int x^n e^x\,dx\)、\(\int x^n\sin x\,dx\) 等积分中反复出现。

反过来，如果选错了 \(u\) 和 \(dv\) 的角色，分部会把你引向死路。对 \(\int x e^x\,dx\)，如果错误地尝试 \(u=e^x\)、\(dv=x\,dx\)，则 \(du=e^x\,dx\)，\(v=x^2/2\)：

\[
\int x e^x\,dx = \frac{x^2}{2}e^x - \int \frac{x^2}{2}e^x\,dx.
\]

新积分比原积分还多了一个 \(x\) 的幂次——更复杂了，不是更简单了。选错方向不丢人，但应该立刻认出来并回头换策略，而不是继续往下多做几层把它包成一个更大的错误。

\subsection{选 \(u\) 和 \(dv\) 的策略：两条经验规则及其边界}

没有万能的选择公式，但两条启发规则覆盖了大多数情形。

第一，选 \(u\) 倾向求导后变简单的东西。幂函数 \(x^n\) 求导降一次——\(x^3\) 变 \(3x^2\)，再变 \(6x\)，再变 \(6\)，最终消失。对数 \(\ln x\) 求导变成 \(1/x\)，从超越降到代数。反三角函数也类似。

第二，选 \(dv\) 倾向容易积出 \(v\) 的东西。\(e^x\)、\(\sin x\)、\(\cos x\) 的原函数就是自己或差个符号——放进 \(dv\) 几乎没代价。

但对 \(\int\ln x\,dx\) 这个看似该让 \(\ln x\) 进入 \(dv\) 的积分——等一下，\(\ln x\) 的原函数不是初等的简洁形式。反过来：令 \(u=\ln x\)，\(dv=dx\)。则 \(du=dx/x\)，\(v=x\)：

\[
\int\ln x\,dx = x\ln x - \int x\cdot\frac{1}{x}\,dx = x\ln x - x + C.
\]

这个例子的教训：选择策略是启发，不是法规。如果首选走不通——新积分比老积分还复杂——立刻回溯，调换 \(u\) 和 \(dv\) 的角色再试。一般不超过两三步就能判断方向对不对。

\subsection{循环积分：分部一次解不出来，代数来凑}

\(\int e^x\cos x\,dx\)。尝试 \(u=\cos x\)，\(dv=e^x\,dx\)，则 \(du=-\sin x\,dx\)，\(v=e^x\)：

\[
I = \int e^x\cos x\,dx = e^x\cos x - \int e^x(-\sin x)\,dx
= e^x\cos x + \int e^x\sin x\,dx.
\]

冒出了 \(\int e^x\sin x\,dx\)——和原来不一样但味道相同。不慌，对这个新积分再分部一次。\(u=\sin x\)，\(dv=e^x\,dx\)：

\[
\int e^x\sin x\,dx = e^x\sin x - \int e^x\cos x\,dx = e^x\sin x - I.
\]

代回第一行：

\[
I = e^x\cos x + (e^x\sin x - I) \;\Longrightarrow\; 2I = e^x(\sin x + \cos x),
\]

\[
I = \frac{e^x}{2}(\sin x + \cos x) + C.
\]

原积分 \(I\) 在等式两边同时出现，这不是循环论证——第二次分部后 \(I\) 的系数是 \(-1\)，代数移项解出 \(I\) 是合法的。每一步都来自乘积法则的逆向使用，不靠直觉跳步。

\subsection{符号技巧有顶，积分的存在不依赖它}

换元和分部覆盖了大量常见积分，但它们有明确的覆盖边界。不是所有连续函数的原函数都能用初等函数（多项式、指数、对数、三角和它们的有限次复合）写出来。最著名的例子是 \(e^{-x^2}\)——处处连续，在任何有限区间上定积分完美存在且唯一，但不定积分没有初等闭合表达式。这个函数以后将以正态分布密度的身份在概率论中反复出现，概率值正是通过数值积分或特殊函数（误差函数 erf）来计算的。

''使用了换元和分部但没积出来''与''这个原函数根本不属于初等函数族''是两个层次的问题。前者可以继续换策略、换方法；后者意味着在初等函数的语言里无解——需要数值积分或引入新的特殊函数。

每一种积分技巧都只匹配特定结构的被积式。换元要求内层与导数因子并存；分部要求被积式可拆成两块且新积分更友好。超出这些结构范围时，强行套用只会越套越乱。知道什么时候换策略、什么时候停手用数值方法——和掌握技巧本身一样是积分功底的一部分。

最后记住一个实用的判断顺序：拿到积分，先看有没有内层加导数因子（换元），再看是不是乘积且其中一项易于求导或降次（分部），两样都不占就看是不是基本积分表覆盖的形状。如果三关全不通过，被积函数连续的话，定积分数值存在——Simpson 法则、Gauss 求积等数值工具随时待命。符号积分和数值积分不是竞争关系，是同一问题的两种开口方式。
